\head{Grades K--1 (7 pages): breaks, last break wins, 2-Sprigs}
\emph{Key in plain words for the adult:} a bar always takes one break fewer than it has squares, in any order, so on a given bar the same seat wins every game. In 2-Sprigs every game uses up every token, two per line, so there are as many lines as dots. Let children find ``it's the same every time''; the reasons are for the adult, and for a child who asks.

\sub{Problem 1, page 1: break three bars}
Bars from the pair's envelope (materials pages 1--2). 3 breaks for the 2\by2, 5 for the 2\by3, 5 for the 1\by6. Children may notice that two different-looking bars take the same number. If two counts for the same bar differ, look for a missed counter or two pieces cut together, and break a bar from the table's spare pool with the adult watching. Prompt before the second bar: ``How many counters do you think will go in the cup?''

\sub{Problem 2, page 2: as few breaks as you can}
11 breaks, whatever the order; nobody can do it with fewer, or with more. The discovery is the table comparison: pairs that broke differently all get 11. A child's explanation: ``every break made one more piece; we started with one piece and ended with 12 squares.'' Hint, for a child sure a clever order saves breaks: ``Which break would you skip? Try it on your second 3~by~4.''

\sub{Problem 3, page 2: last break wins}
3\by3: 8 breaks, the 2nd player wins every game. 2\by4: 7 breaks, the 1st player wins every game. Because the starting player changes, each child wins once on each bar, so the pattern shows in the circles, not in the names. Hint: ``Look at your circles for the square bar. What do you notice?'' Extension: ``Which bar would you want to go first on?''

\sub{Problem 4, page 3: first 2-Sprigs games}
2 dots: 2 lines, 2nd player wins. 3 dots: 3 lines, 1st player wins. Every game. Usual slips: a loop at a dot with one token (third rule picture), a line through a dot (fifth picture), tokens not taken away. A game is over when no line can be drawn. Hint: ``Are there any tokens left? Can anyone still draw?''

\sub{Problem 5, page 4: count the lines}
Every 4-dot game has 4 lines and the 2nd player wins. The tokens check the count: 8 tokens, two per line. Hint: ``How many tokens did each line take away?''

\sub{Problem 6, page 5: can the 2nd player win with 5 dots?}
No. Every game has 5 lines, and the 1st player draws lines 1, 3 and 5. ``It never happened in our four games'' is an honest experimental answer; ``10 tokens, two for each line'' is the explanation a K--1 child can give. Why the game cannot stop with tokens left is in the overview and need not come up. If a pair says the 2nd player won, replay that game together and look for an illegal line.

\sub{Problem 7, page 6: guess, then play}
6 dots: 6 lines, the 2nd player wins both games. Any guess is fine; the goal is a reason such as ``like the 4 dots'' or ``six lines, so the 2nd player draws the last one''.

\sub{Problem 8, page 7: how many breaks made each group?}
2, 4 and 6. The groups have 3, 5 and 7 pieces (for example from a 2\by3, a 3\by3 and a 3\by4 bar), and each break makes one more piece. The tempting answers 5, 8 and 11 are what breaking the whole bar into single squares takes. Check by breaking the spare pool's 2~by~3, 3~by~3 and 3~by~4 (materials page 2) into the same pieces with the cup. For the third group: break off the 1\by4 row, then a 2\by1 column from the 2\by4 left over, split the 2\by3 into two 1\by3 strips, and cut those into a 1\by2 and four single squares with three more breaks, 6 in all. Hint: ``How many pieces were there before the first break? After one break?''

\head{Grades 2--3 (7 pages): predict, explain, then Sprouts}
\sub{Problem 1, page 1: last break wins}
Bars from the pair's envelope (materials pages 1--2). 3\by3: 8 breaks, the 2nd player wins both games. 2\by4: 7 breaks, the 1st player wins both. No: every way of breaking a bar takes the same number of breaks, so the loser cannot change the result. Children first answer from evidence (``we tried different ways''); the explanation often arrives during Problem 2 or 3. Follow-up: ``How do you know a cleverer order would not change it?''

\sub{Problem 2, page 1: fewest breaks}
11, 5 and 3. The fewest is also the most. A child hunting for a shortcut, for example three strips first, can try it on the pair's second 3~by~4 and still counts 11. If a pair reports fewer, look for stacked pieces.

\sub{Problem 3, page 2: decide without breaking}
4\by5: 19 breaks, 1st. 3\by5: 14, 2nd. 5\by5: 24, 2nd. 1\by9: 8, 2nd. 6\by6: 35, 1st. Explanation: breaks $=$ squares $-1$, because every break makes one more piece, and an odd number of breaks means the 1st player makes the last. ``Odd by odd is a 2nd-player bar'' is a correct rule; ask for its reason. Hint: ``How many squares did the 3\by3 have, and how many breaks did it take?'' Extension: which bars would you want to go second on? (Both sides odd.)

\sub{Problem 4, page 3: 2-Sprigs on 3 and 4 dots}
3 dots: 3 moves, 1st player. 4 dots: 4 moves, 2nd player. Every game. The moves record is the evidence Problem 5 needs. A different count means an illegal move (a loop at a one-token dot, a line through a dot or across itself) or a token not taken away.

\sub{Problem 5, page 4: decide first; 10 and 25 dots}
5 dots: 1st (5 moves). 6 dots: 2nd (6 moves). 10 dots: 2nd. 25 dots: 1st. A full Grades 2--3 explanation: ``Each dot has 2 tokens and each line uses 2, so if the game uses all the tokens there are as many lines as dots; an odd number of dots means the 1st player draws the last line.'' At this level ``the game uses all the tokens'' rests on experiment; the follow-up ``Could a game stop while a dot still has a token?'' is Grades 4--5 Problem 5. Hint: ``Look at your moves on page 3. What do you notice?''

\sub{Problem 6, page 5: Sprouts with 1 and 2 spots}
1 spot: always 2 moves, so the 2nd player wins. The first move must be a loop; the old spot and the new spot then have one token each and are joined; the third spot is alone with one token. 2 spots: 4 or 5 moves, so either player can win. Watch for a forgotten new spot or its token, which shortens a game.

\sub{Problem 7, page 6: longest and shortest 2-spot games}
Longest 5, shortest 4; yes, the choices change the winner (5 moves, 1st player; 4 moves, 2nd). With spots A and B and the new spots lettered in order, five moves: A--B (C); A--B again, making a ring A--C--B--D (D); inside the ring A--C (E); still inside, in the part with E, B--D (F); E--F (G), and only G keeps a token. Four moves: A loops (C); inside the loop A--C (D), which shuts D in alone; B loops (E); B--E (F). Hints: for the short game, ``Can you leave a spot with a token where no other spot can reach it?''; for the long game, ``Can every spot that still has a token reach the others?''

\sub{Problem 8, page 7: longest game with 3 spots}
8. There are 9 tokens at the start; each move takes away 2 and adds 1, so after every move there is one token fewer; the newest spot always keeps its token, so at least one is left. An 8-move game, with spots A, B, C and the new spots lettered in order: A--B (D); B--C (E); A--C (F), which closes a ring; then, all on one side of the ring, A--D (G), G--B (H), H--E (I), I--C (J), J--F (K); only K keeps a token. A full answer has both a game of 8 and the reason there is no game of 9. Hint: ``How many tokens are on the board at the start? After one move? After two?'' Extension: 10 spots (29; Grades 4--5 Problem 8).

\head{Grades 4--5 (9 pages): proofs, bounds and Euler's formula}
\sub{Problems 1--2, page 1: last break wins; fewest and most}
Bars as in Grades 2--3, with a spare 3~by~3, 2~by~4 and 3~by~4 for the referee. Problem 1 as in Grades 2--3. Problem 2: fewest $=$ most $=11$, since the pieces go from 1 to 12, one per break. Extension, a rule change: if pieces could be stacked, 4 cuts would do and 3 never could, since a cut at most doubles the number of pieces. This is why the rules say never stack.

\sub{Problem 3, page 2: a rule for any bar}
7\by9: 63 squares, 62 breaks, 2nd. 8\by13: 104 squares, 103 breaks, 1st. 25\by40: 1000 squares, 999 breaks, 1st. Rule: the 1st player wins exactly when the number of squares is even, that is, when a side is even, because there are always squares $-1$ breaks. Starting with a 2\by3 and a 2\by4: 14 squares in 2 pieces, so 12 breaks, and the 2nd player wins. The slip ``$14-1=13$'' misses that the game starts with two pieces; ask ``How many pieces are on the table before the first break?''

\sub{Problem 4, page 3: 2-Sprigs on 3 and 4 dots}
As Grades 2--3 Problem 4: 3 moves and the 1st player; 4 moves and the 2nd.

\sub{Problem 5, page 4: can a game stop with a token left?}
No, so with $n$ dots every game lasts $n$ moves and the 1st player wins exactly when $n$ is odd (proof in the overview). Attempts that fail: a loop from another dot drawn round a one-token dot cannot cross that dot's chain, so it surrounds the whole chain and both ends stay together; two untouched dots in different regions can each still loop. A child's proof: ``A dot with one token is the end of a chain. Follow the chain; its other end has one token too, and nothing can get between them, so they can be joined.'' If a child succeeds, look for a line that crosses itself or a loop at a one-token dot, both forbidden. Hint: ``What do you know about a dot that has only one token left?'' Extension: 3-Sprigs, with three tokens per dot: 4 dots can last 5 or 6 moves, so the winner is no longer fixed.

\sub{Problem 6, page 5: Sprouts, 2 spots}
Play on the two boards of materials page 3 (spots 90\,mm apart, with room for counters and new spots) and write the moves on the page. Longest 5, shortest 4; yes, choices change the winner. Example games are under Grades 2--3 Problem 7. Background: with best play the 2nd player wins with 2 spots.

\sub{Problem 7, page 6: Sprouts, 3 spots}
Longest 8, shortest 6 (7 also occurs). Six: for each spot in turn, loop, then join the spot to the loop's new spot inside the loop, which shuts that new spot in alone. Eight: keep every spot with a token on the edge of one region, joining two that are next to each other along its edge or looping a spot without enclosing another spot with a token; the game under Grades 2--3 Problem 8 is one example. Just never shutting a spot in alone is not enough (see the overview). Hint: ``How many tokens are left at the end of each game?'' (Moves $=9-{}$tokens left.)

\sub{Problem 8, page 7: 10 spots}
29. There are 30 tokens; each move removes two and adds one; the newest spot keeps its token, so at least one remains. The page asks why a game can never last longer; that 29 is actually reached (by the cooperative play of Problem 7, which our check carries out for 10 spots) is a readiness-dependent follow-up. Extension: the fewest is 20, by loop-then-join; proving that nothing shorter is possible is harder (Wikipedia sketches the argument).

\sub{Problem 9, page 7: Brussels Sprouts with 1 and 2 crosses}
1 cross: 3 moves. 2 crosses: 8 moves, however the players try. A different count means a missing bar, a bar with both new arms on one side of its line, a line that crosses itself, or a line touching a cross. Keep the finished drawings for Problem 11.

\sub{Problem 10, page 8: 3 and 10 crosses}
13 and 48. ``Five more for each cross'' from 3 and 8 is a conjecture; the 3-cross game tests it, and $5\times10-2=48$ follows the pattern. Problem 11 gives the reason.

\sub{Problem 11, page 9: dots, lines, regions}
Dots, lines, regions: the tree 5, 4, 1; the square with a diagonal 4, 5, 3; the lens with a loop 2, 3, 3; the finished 1-cross game 4, 6, 4, since each of its three lines passes through the cross it made and counts as two. In every drawing dots $-$ lines $+$ regions $=2$. Finished games from Problem 9: $4-6+4$ and $10-16+8$. For 2 crosses: after $m$ moves there are $2+m$ crosses and $2m$ lines; at the end each region holds exactly one free arm and there are always 8 free arms, so 8 regions; $(2+m)-2m+8=2$ gives $m=8$. Usual slips: 3 lines in the last drawing (giving 5), free arms counted as lines, the outside region forgotten. Follow-ups, to use only when a child is ready: why one free arm per region (at least one, because a line that cuts a region in two sends a bar arm into each half; at most one, because two could still be joined); and why the rule holds for a drawing in one piece (build it a line at a time: a line to a new dot adds a dot and a line, and a line between dots already drawn adds a line and a region). Week 56 proves the formula on solids.

\head{Sources}
Answers above are recomputed by \texttt{check\_guide.py} (in this guide's source folder) against the final pages: every break order for the printed bars; every game of 2-Sprigs with 1--6 dots, Sprouts with 1--3 spots and Brussels Sprouts with 1--2 crosses; region models for 2-Sprigs up to 25 dots and Brussels Sprouts up to 5 crosses; Sprouts games of $3n-1$ and $2n$ moves up to 10 spots, every Sprouts game printed here, and the counterexample in the overview; random Brussels games for the regions-minus-pieces count; and the Problem 11 drawings from their coordinates. The proofs for 2-Sprigs and Brussels Sprouts and the cooperative Sprouts play are written out by us; the $5n-2$ count itself is classical (Gardner, below). The bars and the Grades 4--5 Problem 6 boards come from the theme's materials file, \texttt{materials.pdf} (NN-M-v1, 3 pages); \texttt{check\_guide.py} reads its bars and boards from the PDF to check the counts in the preparation section.

\emph{Consulted on October 25, 2026 (web):} A. Bogomolny, ``Breaking Chocolate Bars'', cut-the-knot (\url{https://www.cut-the-knot.org/proofs/chocolad.shtml}): squares $-1$ breaks, two proofs. JRMF, \emph{Sprigs Teaching Guide} (\url{https://jrmf.org/wp-content/uploads/Sprigs-Teaching-Guide.pdf}), p.~1 (rules: no more than two curve ends at a dot), p.~2 (the 2-Sprigs recording table) and p.~4 (2-Sprigs ends in $n$ moves, stated without proof); its move counts on pp.~4--5 are for 3-Sprigs. ``Sprouts (game)'', Wikipedia (\url{https://en.wikipedia.org/wiki/Sprouts_(game)}): the bounds $2n$ and $3n-1$, and best-play results for 1--7 spots.

\emph{Cited from the outline, not consulted:} D. Fomin, S. Genkin and I. Itenberg, \emph{Mathematical Circles (Russian Experience)}, AMS 1996, and P. Winkler, \emph{Mathematical Puzzles: A Connoisseur's Collection}, 2004 (both cited by cut-the-knot for the bar; no chapter or page given there); M. Gardner, \emph{Mathematical Carnival}, 1975, ch.~1, ``Sprouts and Brussels Sprouts''; E. Berlekamp, J. Conway and R. Guy, \emph{Winning Ways for Your Mathematical Plays}, 2nd ed., vol.~3, 2003, ch.~17, ``Spots and Sprouts''; D. Applegate, G. Jacobson and D. Sleator, 1991, for the remainder-3, 4, 5 conjecture as the outline states it. Library links, also from the outline: Week 56 (Euler's formula on solids), Week 11 (a score that drops every move) and Week 1e (a placement game).
